\chapter{Introduction}
\label{chapter:introduction}

Deriving actionable insights from vast datasets is fundamentally driving innovation \uline{(is a innovation driver across different industries )} across diverse industries, from healthcare analytics to financial modeling. Yet, this heavy reliance on data accumulation \uline{(Yet often times, these innovations stand directly at odds with data privacy, which leads to conflicting...)} presents a critical dichotomy. On one hand, organizations are legally and ethically mandated to protect individual privacy under frameworks such as the \ac{GDPR} and emerging global standards. On the other, they remain wholly dependent on utilizing that same data \uline{Ich würde hier vager bleiben} to train machine learning models and perform complex statistical analyses.

In the specialized field of software engineering, this conflict is centered \uline{Ich wuerde sagen, dass es in der Produktivitätsanalyse von SE darum geht, nicht in SE selbt} on the vast \uline{vast ist ein Wort das die AI sehr oft verwendet, schau mal das du Synonyme findest} amount of metadata generated by development teams through their daily interactions with tools like GitHub and Jira. Effectively harnessing this information to drive productivity while complying with privacy laws requires innovative frameworks that can navigate the complexities of high-dimensional data.

The MECOIS \uline{Ist das ne Abkürzung?} project, a research initiative from the University of Erlangen, addresses the "Big Data" \uline{(this)} challenge within the specific domain of software engineering by transforming the vast metadata generated during development into actionable \uline{auch Wortwiederholung, klingt dadurch sehr AI} intelligence. Software development teams leave a continuous trail of activity-based artifacts in repositories such as GitHub and Jira \uline{erwahenst du auch zweimal identisch, brauchst du nur einmal namedroppen}, creating a high-dimensional data set that is traditionally difficult to interpret \uline{(is difficult to interpret traditionally)} due to overlapping metrics, complexity, and incompatible value ranges. MECOIS solves these issues by applying statistical methods— \uline{AI dash} specifically the creation of a composite index—to group correlating metrics into key dimensions. This data-driven approach \uline{(attepts to provide)} provides an objective Productivity Index that allows organizations to compare performance across teams and projects, identifying early-stage bottlenecks like miscommunication or resource gaps while uncovering the best practices of high-performing teams. \autocite{mecoisProductivity}

To address the tension between data utilization and privacy, this thesis proposes a prototype for a dedicated anonymization service designed to process the metadata collected within the MECOIS framework. While MECOIS extracts high-dimensional, activity-based artifacts from platforms such as GitHub and Jira \uline{schon wieder} to provide deep \uline{deep streichen} insights into team productivity, the granular nature of this data — \uline{AI dash} including timestamps, commit details, and developer identities — presents significant privacy risks. The proposed service aims to anonymize this data in strict accordance with the \ac{GDPR} \uline{Hier vielleicht noch nicht namedroppen sodern einfach sagen mit Gesetz und ethischen Massstäben}, ensuring that the trail of artifacts used to calculate the MECOIS Productivity Index is stripped of personally identifiable information. By integrating this anonymization layer, the project seeks to enable the MECOIS’s mission of empowering software developers through data-driven analysis while maintaining a privacy-first approach that meets modern legal and ethical standards.